\section{Formalizacja problemu RL/MDP}
\label{sec:formalizacja}

Problem modelujemy jako epizodyczny proces decyzyjny Markowa
$(\mathcal{S}, \mathcal{A}, P, R, \gamma)$.

\subsection{Agent i środowisko}

\begin{itemize}
  \item \textbf{Agent} -- dron; w każdym kroku wybiera jedną akcję.
  \item \textbf{Środowisko} -- mapa $W \times H$ wraz z bazą, celami,
        stacjami ładowania, przeszkodami, strefami radarowymi, patrolem
        i korytarzem zamykanym. Środowisko przechowuje bieżący stan, wykonuje
        przejścia i zwraca nagrody.
\end{itemize}

\subsection{Stan lub obserwacja}

Stan środowiska to krotka
\begin{equation}
  s = (x,\; y,\; e,\; m,\; \varphi),
\end{equation}
której składowe opisano w tabeli~\ref{tab:stan}.

\begin{table}[h]
\centering
\begin{tabular}{p{0.30\textwidth}p{0.60\textwidth}}
\toprule
Składowa stanu/obserwacji & Znaczenie \\
\midrule
$(x, y)$ & pozycja drona na mapie \\
$e \in \{0, 1, \dots, E_{\max}\}$ & poziom baterii \\
$m \in \{0,1\}^K$ & maska bitowa sfotografowanych celów
  ($m_i = 1$ -- cel $i$ został sfotografowany) \\
$\varphi \in \{0, \dots, T_p - 1\}$ & faza patrolu, $\varphi = t \bmod T_p$,
  gdzie $T_p$ to okres trasy patrolu \\
\bottomrule
\end{tabular}
\caption{Składowe stanu.}
\label{tab:stan}
\end{table}

\paragraph{Stan a obserwacja.} Agent obserwuje pełny stan $s$, więc problem
jest w pełni obserwowalnym MDP. Agent nie dostaje natomiast informacji
o modelu środowiska: nie zna układu mapy, trasy patrolu, wartości
$p_{\text{det}}$ ani skutków sfotografowania celu~1. Elementy te należą do
dynamiki przejść $P$, a nie do stanu, i agent poznaje je wyłącznie pośrednio,
przez obserwowane przejścia i nagrody. Pozycja patrolu nie jest podawana
wprost, ponieważ wynika jednoznacznie z fazy $\varphi$.

\paragraph{Dlaczego faza patrolu jest w stanie?} Pozycja patrolu zależy od
czasu. Gdyby $\varphi$ nie było częścią stanu, ten sam stan $(x,y,e,m)$
prowadziłby do różnych skutków w zależności od chwili, w której agent się
w nim znalazł -- własność Markowa nie byłaby spełniona. Pozycję patrolu da się
odtworzyć z $\varphi$, więc nie trzeba jej przechowywać osobno.

\paragraph{Dlaczego stan korytarza $Z$ nie jest w stanie?} Korytarz jest
zamknięty dokładnie wtedy, gdy sfotografowano cel~1, czyli gdy $m_1 = 1$.
Informacja ta jest już zawarta w masce $m$. Gdyby zdarzenie zależało od innego
warunku (np. czasu), stan należałoby rozszerzyć o dodatkową flagę.

\paragraph{Rozmiar przestrzeni stanów.} Górne oszacowanie wynosi
\begin{equation}
  |\mathcal{S}| \le W \cdot H \cdot (E_{\max}+1) \cdot 2^K \cdot T_p .
\end{equation}
Dla mapy \emph{easy} ($10\times10$, $E_{\max}=30$, $K=3$, $T_p=6$) daje to
$100 \cdot 31 \cdot 8 \cdot 6 = 148\,800$ stanów, a po uwzględnieniu
$|\mathcal{A}| = 5$ około $7{,}4 \cdot 10^5$ wartości $Q(s,a)$. Taka tablica
mieści się w pamięci (kilka MB), więc podejście tablicowe jest wykonalne.
Faktycznie osiągalnych stanów jest mniej, ponieważ pola z przeszkodami
nie są dostępne.

\subsection{Akcje}

\[
  \mathcal{A} = \{\textsc{góra},\ \textsc{dół},\ \textsc{lewo},\ \textsc{prawo},\ \textsc{czekaj}\}.
\]

\begin{table}[h]
\centering
\begin{tabular}{p{0.30\textwidth}p{0.60\textwidth}}
\toprule
Akcja & Opis \\
\midrule
\textsc{góra} & ruch o jedno pole w górę \\
\textsc{dół} & ruch o jedno pole w dół \\
\textsc{lewo} & ruch o jedno pole w lewo \\
\textsc{prawo} & ruch o jedno pole w prawo \\
\textsc{czekaj} & pozostanie na miejscu; na stacji ładowania zwiększa energię
  o $\Delta_L$ (do $E_{\max}$), poza stacją oznacza zawis i zużywa energię
  jak ruch; pozwala przeczekać przejście patrolu \\
\bottomrule
\end{tabular}
\caption{Zbiór akcji agenta.}
\label{tab:akcje}
\end{table}

Wszystkie akcje są dostępne w każdym stanie, dzięki czemu zbiór
$\mathcal{A}$ jest stały, a tablica $Q$ ma prostą strukturę. Ruch na pole
poza mapą, na przeszkodę lub na zamknięty korytarz jest \textbf{akcją
niepoprawną}: dron pozostaje na miejscu, zużywa energię i otrzymuje
dodatkową karę.

\subsection{Przejście (pojedynczy krok interakcji)}

W kroku $t$ środowisko wykonuje kolejno:
\begin{enumerate}
  \item wyznaczenie nowej pozycji na podstawie akcji $a_t$ (z obsługą akcji
        niepoprawnych),
  \item aktualizację energii: $e \leftarrow e - 1$ albo
        $e \leftarrow \min(e + \Delta_L, E_{\max})$ przy ładowaniu,
  \item przesunięcie patrolu: $\varphi \leftarrow (\varphi + 1) \bmod T_p$,
  \item sprawdzenie kolizji z patrolem (również „minięcia się'' na sąsiednich
        polach, gdy dron i patrol zamieniają się miejscami),
  \item jeżeli dron jest na niesfotografowanym celu $i$ -- ustawienie $m_i = 1$
        (dla $i=1$ zamknięcie korytarza $Z$),
  \item jeżeli dron jest w strefie radarowej -- losowanie wykrycia
        z prawdopodobieństwem $p_{\text{det}}$,
  \item sprawdzenie warunków zakończenia i zwrócenie $(s_{t+1}, r_{t+1},
        \textit{terminated}, \textit{truncated})$.
\end{enumerate}

Interfejs środowiska będzie zgodny z biblioteką Gymnasium~\cite{gymnasium}
(\texttt{reset()}, \texttt{step(a)}).

\subsection{Funkcja nagrody}
\label{sec:nagroda}

Nagroda otrzymana po wykonaniu akcji $a_t$ w stanie $s_t$ wynosi
\begin{equation}
  r_{t+1} = -1
    - 2 \cdot \mathbf{1}[\text{akcja niepoprawna}]
    + 10 \cdot \mathbf{1}[\text{nowy cel}]
    + 50 \cdot \mathbf{1}[\text{sukces}]
    - 50 \cdot \mathbf{1}[\text{porażka}],
\end{equation}
gdzie $\mathbf{1}[\cdot]$ jest funkcją wskaźnikową, a zdarzenia
i uzasadnienie wartości zebrano w tabeli~\ref{tab:reward}.

\begin{table}[h]
\centering
\begin{tabularx}{\textwidth}{>{\raggedright}p{0.36\textwidth} r >{\raggedright\arraybackslash}X}
  \toprule
  Zdarzenie & Nagroda & Uzasadnienie \\
  \midrule
  każdy krok & $-1$ & zachęca do krótkich tras \\
  akcja niepoprawna & $-2$ (dodatkowo) & zniechęca do „odbijania się'' od ścian \\
  sfotografowanie nowego celu & $+10$ & nagroda pośrednia, ułatwia eksplorację \\
  powrót do bazy ze wszystkimi celami & $+50$ & główny cel misji \\
  wykrycie, kolizja z patrolem & $-50$ & porażka misji \\
  rozładowanie baterii & $-50$ & porażka misji \\
  \bottomrule
\end{tabularx}
\caption{Wstępna funkcja nagrody (wartości do strojenia po pierwszych testach).}
\label{tab:reward}
\end{table}

\paragraph{Kara za krok.} Kara $-1$ za każdy krok, także za ładowanie
i zawis, sprawia, że krótsza trasa daje wyższy zwrot. Udana misja na mapie
\emph{easy} (ok.~50 kroków razem z ładowaniem) daje zwrot około
$-50 + 3 \cdot 10 + 50 = +30$, czyli wyraźnie więcej niż dowolna porażka.

\paragraph{Ograniczenia twarde i miękkie.} Ograniczenia \emph{twarde} są
narzucone przez środowisko niezależnie od nagrody: granice mapy, przeszkody
i zamknięty korytarz $Z$ (ruch jest niemożliwy) oraz rozładowanie baterii,
kolizja z patrolem i limit kroków (koniec epizodu). Ograniczeniem
\emph{miękkim} są strefy radarowe: wejście do nich jest dozwolone, ale wiąże
się z ryzykiem wykrycia, a o tym, czy skrót się opłaca, agent decyduje na
podstawie oczekiwanego zwrotu. Miękka jest także kara $-2$ za akcję
niepoprawną -- zniechęca do niej, ale jej nie zabrania.

\paragraph{Czy nagroda nie zachęca do niepożądanego zachowania?}
\begin{itemize}
  \item Nagroda $+10$ jest przyznawana tylko za \emph{nowy} cel
        ($m_i$ zmienia się z $0$ na $1$), więc nie da się jej zbierać
        wielokrotnie, krążąc wokół celu.
  \item Ładowanie nie jest nagradzane, więc agent nie ma powodu ładować się
        ponad potrzebę.
  \item Zebranie celów bez powrotu kończy się porażką lub obcięciem epizodu,
        więc powrót do bazy zawsze się opłaca.
  \item \textbf{Ryzyko: porażka może być atrakcyjniejsza niż długie
        przetrwanie.} Agent, który nie potrafi jeszcze ukończyć misji
        i krąży (np. ładując się na stacji) aż do limitu $T_{\max} = 200$,
        uzyskuje zwrot około $-200$, natomiast szybkie wejście w strefę
        radarową kończy się zwrotem rzędu $-55$. Suma kar za kroki może więc
        przewyższyć karę za porażkę i na wczesnym etapie uczenia agent może
        nauczyć się celowego kończenia epizodu. Możliwe przeciwdziałania to
        zmniejszenie kary za krok (np. do $-0{,}1$), zwiększenie kary za
        porażkę albo skrócenie $T_{\max}$. Wybór zostanie dokonany po
        pierwszych testach (zob. sekcja~\ref{sec:otwarte}).
\end{itemize}

\subsection{Warunki zakończenia epizodu}
\label{sec:koniec}

\begin{itemize}
  \item \textbf{Sukces} (\textit{terminated}): dron jest w bazie i $m = (1,\dots,1)$.
  \item \textbf{Porażka} (\textit{terminated}): wykrycie w strefie radarowej,
        kolizja z patrolem lub $e = 0$ poza bazą.
  \item \textbf{Przekroczenie limitu} (\textit{truncated}): $t \ge T_{\max}$.
\end{itemize}

Rozróżnienie \textit{terminated}/\textit{truncated} jest istotne przy
aktualizacji wartości: po zakończeniu właściwym wartość stanu następnego
wynosi $0$, natomiast po obcięciu limitem kroków agent nadal powinien
korzystać z oszacowania $\max_a Q(s_{t+1},a)$.

\subsection{Założenia i ograniczenia środowiska}

\paragraph{Uproszczenia modelu.}
\begin{itemize}
  \item Przestrzeń jest dyskretna (siatka), dron porusza się w czterech
        kierunkach o jedno pole; pominięto wysokość lotu, prędkość
        i bezwładność.
  \item Każdy krok trwa tyle samo i kosztuje tyle samo energii, niezależnie
        od kierunku oraz od tego, czy dron leci, czy wisi w miejscu.
  \item Ruch jest deterministyczny -- nie modelujemy wiatru ani błędów
        sterowania (wiatr jest rozważany jako rozszerzenie).
  \item Zdjęcie jest wykonywane automatycznie i natychmiast po wejściu na pole
        celu; pominięto orientację kamery i warunki widoczności.
  \item Wykrycie w strefie radarowej to niezależne losowanie w każdym kroku
        ze stałym prawdopodobieństwem $p_{\text{det}}$, bez zależności od
        odległości od radaru ani czasu przebywania w strefie.
  \item Występuje jeden patrol poruszający się po stałej, cyklicznej trasie,
        który nie reaguje na drona.
  \item Mapa jest statyczna z wyjątkiem korytarza $Z$; stacje ładowania mają
        nieograniczoną pojemność i stałe tempo ładowania.
  \item Agent obserwuje pełny stan -- pominięto szum czujników i ograniczone
        pole widzenia.
\end{itemize}

\paragraph{Parametry środowiska.}
\begin{itemize}
  \item energia $E_{\max}$ (wstępnie 30) i przyrost ładowania $\Delta_L$
        (wstępnie 5 na krok),
  \item limit kroków $T_{\max}$ (wstępnie 200),
  \item prawdopodobieństwo wykrycia $p_{\text{det}}$ (wstępnie 0{,}3 na krok),
  \item brak możliwości wejścia na przeszkody i zamknięty korytarz.
\end{itemize}

\subsection{Współczynnik dyskontowania}

Wstępnie $\gamma = 0{,}99$. Główna nagroda pojawia się dopiero na końcu
długiego epizodu, więc zbyt mała wartość $\gamma$ sprawiłaby, że agent
praktycznie nie odczuwałby nagrody za powrót do bazy. Wpływ $\gamma$ zostanie
sprawdzony w eksperymentach.

\subsection{Miary sukcesu projektu}
\label{sec:miary}

\begin{itemize}
  \item odsetek misji zakończonych sukcesem (w fazie ewaluacji, bez eksploracji),
  \item średni zwrot z epizodu,
  \item średnia długość udanego epizodu,
  \item odsetek porażek z podziałem na przyczyny (wykrycie, patrol, bateria),
  \item stabilność wyników między uruchomieniami z różnym ziarnem losowym.
\end{itemize}

Projekt uznajemy za udany, jeżeli agent uczący się istotnie przewyższa agentów
bazowych na wszystkich mapach, a zachowanie agenta (np. omijanie radarów,
planowanie ładowania, kolejność celów) daje się zinterpretować.

\subsection{Elementy wymagające weryfikacji po pierwszych testach}
\label{sec:otwarte}

\begin{itemize}
  \item wartości nagród -- czy agent nie uczy się „samobójczych'' polityk
        albo krążenia w miejscu; w szczególności bilans kary za krok i kary
        za porażkę (sekcja~\ref{sec:nagroda}),
  \item czy nagroda $+10$ za cel nie powoduje, że agent zbiera cele i nie
        wraca do bazy,
  \item wartość $p_{\text{det}}$ -- zbyt duża sprawi, że radary staną się
        zwykłymi ścianami, zbyt mała -- że będą bez znaczenia,
  \item czy osiągalna przestrzeń stanów jest wystarczająco mała dla
        tablicowego Q-learningu na mapie \emph{hard},
  \item ewentualne dodanie losowego znoszenia drona przez wiatr
        (dodatkowe źródło niepewności).
\end{itemize}
